% Proposta de Projeto — Aprendizado de Máquina (UFPR)
% Baseado no template IEEE/CVPR

\documentclass[10pt,twocolumn,letterpaper]{article}

%%%%%%%%% PAPER TYPE
\usepackage{cvpr}

% Pacotes
\usepackage{graphicx}
\usepackage{amsmath}
\usepackage{amssymb}
\usepackage{booktabs}
\usepackage[utf8]{inputenc}
\usepackage[T1]{fontenc}
\usepackage[brazil]{babel}
\usepackage[pagebackref,breaklinks,colorlinks]{hyperref}
\usepackage[capitalize]{cleveref}
\crefname{section}{Seção}{Seções}
\Crefname{section}{Seção}{Seções}
\Crefname{table}{Tabela}{Tabelas}
\crefname{table}{Tab.}{Tabs.}

\def\cvprPaperID{*****}
\def\confName{CVPR}
\def\confYear{2022}

\begin{document}

%%%%%%%%% TÍTULO
\title{Proposta de trabalho}

\author{
Pedro Henrique Marques de Lima \\
Universidade Federal do Paraná \\
{\tt\small pedro.marques@ufpr.br}
\and
Felipe Pereira Gonçalves \\
Universidade Federal do Paraná \\
{\tt\small felipe.goncalves@ufpr.br}
}
\maketitle

% =====================================================================
\section{Problemas e Aplicações}

Enchentes e deslizamentos provocados por chuvas extremas causam
centenas de mortes e bilhões de reais em prejuízos no
Brasil~\cite{marengo2020trends,cenad2024}. Antecipar esses eventos a
partir de variáveis meteorológicas observadas é relevante para a
proteção civil.

Deepfakes são conteúdos realistas gerados por redes generativas, o problema ocorre no uso malicioso para aplicação de golpes, disseminação de desinformação e chantagem usando imagem rostos de pessoas reais~\cite{Rafique2023,Mary2023}.

\textbf{Problema~1 (tabular).}
Classificação binária horária: dado um vetor de variáveis
meteorológicas $\mathbf{x}\!\in\!\mathbb{R}^{n}$, prever se o
acumulado de precipitação excede um limiar crítico ($y\!=\!1$).
Os desafios são o desbalanceamento de classes (chuvas extremas são
raras), dados faltantes nas séries do INMET e a dependência temporal
entre observações.


\textbf{Problema~2 (não tabular).} Classificação binária: dada uma imagem em formato jpeg determinar se a imagem foi gerada artificialmente ($y\!=\!1$).
 O principal desafio esta na sofisticação das manipulações, que as tornam quase indistinguíveis ao olho humano.

% =====================================================================
\section{Conjunto de Dados}

Os dados do problema 1 provêm do INMET\footnote{\url{https://portal.inmet.gov.br/dadoshistoricos}}
(observações horárias de estações automáticas, 2022--2025). Foram
selecionados cinco eventos reais de precipitação extrema
(\cref{tab:eventos}); para cada um, a estação mais próxima foi
escolhida entre 150 disponíveis nos estados alvo (PR, SP, RS, PE),
com recorte de $\pm$4~meses em torno do desastre. 

\begin{table}[h]
\centering
\caption{Eventos selecionados e estações INMET.}
\label{tab:eventos}
\small
\setlength{\tabcolsep}{3pt}
\begin{tabular}{@{}llll@{}}
\toprule
\textbf{Região} & \textbf{UF} & \textbf{Estação} & \textbf{Data} \\
\midrule
Serra do Mar / Guaratuba  & PR & Morretes (A873)   & 28/11/2022 \\
Litoral Norte             & SP & Bertioga (A765)   & 19/02/2023 \\
Vale do Taquari           & RS & Teutônia (A882)   & 04/09/2023 \\
Bacia do Guaíba           & RS & P.\ Alegre (A801) & 02/05/2024 \\
Região Metropolitana      & PE & Palmares (A357)   & 06/02/2025 \\
\bottomrule
\end{tabular}
\end{table}

O conjunto consolidado possui \textbf{16\,013 registros horários} com dados sobre: precipitação, pressão atmosférica, temperaturas, umidade relativa e vento.

Sobre o problema 2, a base de dados foi obtida em um repositório do Kaggle\footnote{\url{https://www.kaggle.com/datasets/prithivsakthiur/deepfake-vs-real-20k}}, os dados dessa base foram coletados da internet aberta e processados utilizando os pipelines do DeepfakeNet para obter apenas imagens deepfakes de alta qualidade. Os dados são divididos em duas classes balanceadas ideal para uso de acurácia como métrica\cite{Rafique2023}. A base conta com imagem formato jpeg de resoluções variadas sendo 9576 imagens deepfake e 9643 imagens reais.
% =====================================================================
\section{Metodologia}

\subsection{Pré-processamento}
\textbf{Problema 1:} Falhas serão tratadas por interpolação temporal Variáveis derivadas de janelas deslizantes (acumulado 3--12\,h, variação de
pressão) serão adicionadas. Normalização por z-score.

\textbf{Problema 2:} As imagens devem ser redimencionadas para as dimensões exigidas pela rede escolhida. Além do uso de Error Level Analysis (ELA), as imagens serão ressalvas com uma taxa de erro de 95\% para calcular a diferença entre os níveis de compressão. Isso ajuda a destacar áreas manipuladas digitalmente\cite{Rafique2023}.

\subsection{Modelos}
\textbf{Problema 1:}
Random Forest, XGBoost/LightGBM.
Desbalanceamento tratado com SMOTE e ajuste de
pesos.

\textbf{Problema 2:}
SqueezeNet, ResNet18, uso de CNN em conjunto com SVM ou KNN\cite{Mary2023,Rafique2023}.

\subsection{Avaliaçao}
Para ambos problemas, utilizaremos métricas: F1-Score, AUC-ROC, Precision-Recall AUC e Balanced
Accuracy, com matrizes de confusão e curvas ROC/PR. Para o problema 1, a divisão será feita respeitando ordem temporal.

% =====================================================================
\newpage
{\small
\bibliographystyle{ieee_fullname}
\bibliography{egbib}
}

\end{document}
